We shall apply these definitions to the augmented @@M0@@-rings @@M1@@ and @@M2@@, the augmentation rings @@M3@@ and @@M4@@ being binomial (V 2.8). To simplify notation, we shall put, for every scheme @@M5@@ and every @@M6@@:
@@M7@@
For every graded algebra @@M8@@, we shall regard the group @@M9@@ as contained in @@M10@@. Finally, given a homomorphism @@M11@@ which is additive and compatible with the gradings up to a shift, we shall still denote by @@M12@@ the homomorphism @@M13@@ obtained from it.

\sgaStableTarget{sga6.e08.para.3-4}\textbf{3.4.} Let @@M14@@ be a scheme. For every locally free @@M15@@-Module of finite type @@M16@@, denote by @@M17@@ the class in @@M18@@ of the dual @@M19@@ of @@M20@@. It is clear that the function associating @@M22@@ to @@M21@@ is additive on the category of locally free @@M23@@-Modules of finite type, and therefore defines an endomorphism of @@M24@@, which we shall denote by @@M25@@.

\sgaStableTarget{sga6.e08.prop.3-5}\textbf{Proposition 3.5.} \uline{Let} @@M26@@ \uline{be a quasi-compact scheme}. \uline{Then for every} @@M27@@ \uline{and every} @@M28@@, \uline{we have}:
@@M29@@

The relation to be proved may also be written:
@@M30@@
where we put: @@M31@@.

By additivity, one may therefore suppose that @@M32@@ is the class of a locally free @@M33@@-Module of finite type @@M34@@. Moreover, after replacing @@M35@@ by the flag scheme of @@M36@@, which gives an injection on the
